\section{稀疏图码与消息传递}
\label{sec:04-Discrete-Mathematics/05-Coding-and-Polynomial-Methods/04}

手机连着的这个 WiFi 链路，每一帧都在做译码，几百微秒一次，功耗以毫瓦计。它跑的码长上万比特，合法码字的数量是二的几千次方——若按定义去找与接收波形最相容的那个码字，把宇宙拆了当计算机也算不完。

最大似然译码的公式写起来毫不费力：在所有合法 \(c\) 中最大化 \(P(y\mid c)\)。这个式子定义正确，却不构成一个方法。上一节的代数译码给出了一条出路，条件是整个码字要被一个全局的多项式关系拴住，域一大运算就贵，而且信道给出的置信度用不上。

低密度奇偶校验码（LDPC code）走了相反的方向：不用一个刚性的全局约束，改用成百上千条各自只看几位的局部约束。单条约束弱得可笑，条与条之间却互相搭边，于是一次全局搜索被改写成邻居之间反复交换意见。

\subsection{稀疏校验矩阵就是一张二分图}

LDPC 码仍是线性码，由校验矩阵定义，合法码字满足 \(Hc^{\trans}=0\)。特别之处只在于 \(H\) 极其稀疏：每行只涉及少数几个变量，每列也只出现在少数几个方程里，且这个``少数''不随码长增长。

稀疏矩阵画成图更好读。每列一个变量节点（一个码字符号），每行一个校验节点（一条奇偶约束），\(H_{ji}=1\) 就在二者之间连一条边，得到的二分图叫 Tanner 图（二分图的一般性质见\hyperref[ch:028]{``图论：局部邻接怎样形成全局网络''}）。

\begin{center}
\begin{tikzpicture}[x=1cm,y=1cm,font=\small,>=stealth]
  \foreach \a/\b in {0.5/1.7, 1.7/1.7, 4.1/1.7, 1.7/3.5, 2.9/3.5, 5.3/3.5, 2.9/5.3, 5.3/5.3, 6.5/5.3}{
    \draw[black!40] (\a,0.22) -- (\b,2.18);
  }
  \foreach \x/\n in {0.5/1,1.7/2,2.9/3,4.1/4,5.3/5,6.5/6}{
    \draw[black!60,fill=blue!12] (\x,0) circle (0.22);
    \node[font=\footnotesize] at (\x,0) {\(c_\n\)};
    \draw[->,black!40] (\x,-0.9) -- (\x,-0.3);
  }
  \foreach \x in {1.7,3.5,5.3}{
    \draw[black!60,fill=black!10] (\x-0.22,2.18) rectangle (\x+0.22,2.62);
    \node[font=\footnotesize] at (\x,2.4) {\(+\)};
  }
  \draw[->,thick,blue!70] (1.48,0.36) -- (1.48,2.04);
  \node[left,font=\footnotesize] at (1.4,1.25) {我倾向于};
  \draw[->,thick,red!65] (3.66,2.0) -- (3.06,0.38);
  \node[right,font=\footnotesize] at (3.74,1.25) {那你该是};
  \node[below,font=\footnotesize,black!70] at (3.5,-1.0) {信道证据};
  \node[right,font=\footnotesize,black!70] at (6.85,2.4) {校验节点};
  \node[right,font=\footnotesize,black!70] at (6.85,0) {变量节点};
\end{tikzpicture}

\emph{每个变量只认识一两个校验，每个校验也只管三个变量，没有谁看得到全局。译码是这张图上的来回通信：变量向上报告自己目前的倾向，校验向下回复``要让我这条约束成立，你更可能取什么''。轮数一多，本来只在局部的证据就沿着边渗到了远处。}
\end{center}

图不是矩阵的装饰画法。它标出了算法可以利用的结构：每个节点只与邻居打交道，计算量随码长线性增长，而不是随码字数指数增长。

\subsection{擦除信道上，消息传递退化成连锁剥离}

先看一个不需要概率的版本。信道只会擦除，接收方知道哪些位置丢了。若某条校验
\[
c_1+c_2+c_3=0
\]
里 \(c_1,c_2\) 已知而 \(c_3\) 被擦除，这条约束一步就把 \(c_3\) 定死。新恢复的 \(c_3\) 又出现在别的校验里，可能使那条校验也只剩一个未知量，于是解一个、触发一片，像推倒的骨牌。

整个译码没有任何全局计算，只是反复扫描``哪条校验只剩一个未知''。它会不会走到底，取决于图的连接方式：若剩余的未知变量恰好构成一个每条校验都至少碰到两个未知量的子结构，骨牌就推不下去了——这时信息未必不足，只是这套局部规则用尽了。此处的停顿与信息不足是两回事，下一节还要专门算这笔账。

\subsection{软消息，以及不把原话奉还的规矩}

真实信道给的不是比特，是带强弱的证据。收到 \(+0.1\) 与 \(+5.0\)，硬判都算 0，可信度天差地别。一旦在入口处压成 \(0/1\)，那句``前者很可疑''就永久丢失了，后面的迭代再也补不回来。

所以消息用对数似然比携带：
\[
L_i=\log\frac{P(c_i=0\mid y_i)}{P(c_i=1\mid y_i)},
\]
符号表示偏向哪一边，绝对值表示信心多强。变量节点把信道证据与各邻居的意见加起来，校验节点则根据其余邻居的意见推断目标变量该取什么。

这里有一条容易忽略却必须守住的规矩：变量发给某个校验的消息里，不许包含这个校验刚发来的内容；反过来也一样。这叫传递外在信息。若把对方的话原样放大再还回去，两个节点会互相引用对方给出的证据，几轮之内就能把一个毫无根据的判断顶成``高置信''。用统计的话说，这是把同一份数据重复计入了似然。

\subsection{同一个算法，在三个领域各有名字}

若 Tanner 图是一棵树，消息传递给出的边缘概率是精确的。理由与动态规划在树上有效完全相同：切断一条边，两侧条件独立，一份证据不可能沿两条路径重复到达同一个节点。

高性能的码不可能是树——树的叶子太多、码率上不去。Tanner 图必然含环，于是同一份证据绕一圈后以``新消息''的面目回来，独立性假设破裂，迭代结果只是近似。环长的时候问题不大，前几轮里局部邻域看起来仍像一棵树；短环则让证据过早地自我确认。

值得停一下的是，这套``节点只与邻居交换消息、迭代若干轮''的算法，在概率图模型里叫置信传播，用来在因子图上做近似推断，精确条件同样是无环，失效方式同样是短环上的重复计数（见\hyperref[sec:13-Machine-Learning/07-Graphical-and-Sequential-Models/01]{``图结构把联合分布拆成局部因子''}）。校验节点在这里扮演的就是因子的角色，奇偶约束是一个取值为 0 或 1 的因子。图神经网络的邻居聚合也是同一副骨架：每一层把邻居的表示汇总成新的节点表示，堆若干层，信息传播若干跳；层数太多时不同节点的表示趋同，与短环上证据反复回流是同一类病。三个领域各自发展出的记号很不一样，共享的结构是：约束是局部的，推断靠迭代，正确性只在树上有保证。

Turbo 码可以放进同一个框架看。它用交织器把同一批信息位打乱后交给第二个卷积码，两个译码器轮流交换外在信息——一个译码器啃不动的位，在另一种排列下可能有强证据。图结构与 LDPC 不同，做的事情是一样的：把全局约束拆散，在局部算软证据，只交换对方还不知道的部分，指望迭代收敛到一个全局自洽的解释。它的性能来自交织、软输入软输出与迭代三者的合力，没有哪一个单独起决定作用。

\subsection{接近容量是有条件的说法}

LDPC 与 Turbo 码能在许多信道上逼近 Shannon 极限，这句话的适用范围需要说清楚。它是渐近陈述，码长几百的短码通常显不出优势；译码器要有一个像样的信道概率模型，模型偏了软信息就是误导；迭代次数、量化位宽、硬件并行度都会改变实际曲线。

更要紧的是，置信传播在含环图上不提供最大似然保证。误码率曲线常在低噪声区突然变平，形成误码平台，通常由陷阱集造成——一小撮变量彼此支持着一个错误的判断，迭代进去就出不来，尽管全局上存在更优解释。所以 LDPC 的设计不是随手往矩阵里撒 1，行列度分布、短环数量、图的扩张性都要算过。

这正是理论能力与可计算实现的分工：最小距离和信道容量划出边界，图结构与算法决定能否在现实成本里贴近它。

下一节换一个信道模型。若数据包要么完整到达、要么明确丢失，而发送方事先不知道每个接收方会丢多少——比如同一份模型权重要推给上万台机器，各家丢的包各不相同——那么在发送前就把码率定死这件事本身，就成了麻烦的来源。
